\chapter{Execution paths}\label{ch:executing}

The same native program runs on two paths (\cref{fig:bypass}).

\section{Execution through Metal}\label{ch:through-metal}

The executor, \code{tools/g17decodegen.m}, loads each bundle's library and
binary archive. It builds the pipeline with \code{MTLPipelineOptionFailOnBinaryArchiveMiss}, so Metal must use the
archived code. Patching one constant in GPU memory after the pipeline is built changes the output as predicted
(MM~\mm{25.147}), so the GPU runs these bytes; whether they are copied or checked again before execution is not
determined. Metal still derives the pipeline configuration from the metadata the image declares and submits the work.

The executor creates one Metal buffer per arena, writes the arenas' initial contents, and encodes all of a token's
dispatches into one command buffer. Buffers, bindings and dispatches are fixed before the first token, and each
stage's output stays on the GPU in an arena that the next stage binds. In pipelined mode the GPU also runs the generation step, which advances
the position and writes the next token's embedding. The host then commits the token command buffers back to back
with no wait between them, and reads the generated tokens from a log at the end. With no idle gap between tokens the
GPU stays at its maximum clock state (MM~\mm{0.15}).

In the current InternLM2.5-1.8B decode graph each of the 24 layers is seven dispatches:

\begin{lstlisting}
RMSNorm -> qkv projection -> attention (RoPE + KV append + 32 key slices + merge) -> wo projection + residual
        -> RMSNorm -> w1/w3 projections + SwiGLU -> w2 projection + residual
\end{lstlisting}

Four more dispatches run per token: the final norm, the vocabulary head, the argmax and the generation step. That
makes 172 dispatches a token (MM~\mm{25.154}, MM~\mm{25.190}). mlx-lm runs 16 dispatches a layer for the same model
(MM~\mm{25.142.1}). Qwen3-0.6B and Qwen3-8B also run on this path (MM~\mm{25.188}, MM~\mm{25.204}).

Split-K quantized matrix-vector products read the 4-bit
weights with K split across lanes and threadgroups, and fuse the residual add or SwiGLU into the same dispatch
(MM~\mm{25.152}, MM~\mm{25.153}). Attention splits the keys into 32 slices, each a simdgroup with its own running
softmax, and merges the partial results with a parallel butterfly. Attention for one layer at 128 keys takes
9.7~µs, against 12.0~µs in MLX (MM~\mm{25.140}, MM~\mm{25.141.14}).

Metal versions of these kernels, compiled by Apple with identical threads, partitions and operation order, are
bit-identical and 10–15 percent faster (\cref{sec:bounds-real-workloads}); MM~\mm{25.211} has not yet found which
code difference costs the time.

In batched decode, with up to 16 sequences, the projection dequantizes each 4-bit weight once in registers, then multiplies it against every sequence with tensor MMAs. It computes the
transpose, \code{Y\textasciicircum{}T = W\textasciicircum{}T X\textasciicircum{}T}, because only B can advance through the key-block index register (MM~\mm{25.166}). The dequantization runs in fp16 on the register-form fp16 fma, \opn{798} (MM~\mm{25.196}).
Splitting K to about 512 threadgroups raises 16-sequence decode from 1,172 to 1,462 tok/s (MM~\mm{25.185}), and a 32-sequence form runs two 16-row
halves over one dequantized weight (MM~\mm{25.178}).

Prefill runs the prompt in chunks of up to 512 rows. Its GEMMs run on the tensor unit against persistent fp16 weights.
Attention keeps its score and probability tiles in registers, and the FFN intermediate is fp16. Short prompts use a
16-row 4-bit route instead (MM~\mm{25.163}, MM~\mm{25.183}, MM~\mm{25.202}).

\section{Native execution below Metal}\label{ch:below-metal-native}

The native runtime runs a model with no Metal in the process, and it checks that AGXMetal, Apple's Metal driver
library, is absent (\cref{sec:opening-device-contexts}). Work is submitted through Apple's private
\code{IOGPU} framework to the installed kernel driver and the GPU firmware (\cref{tab:native-path-ownership}). Two models run on it, MiniLM-L6
(sentence embeddings) and Qwen2.5-0.5B-Instruct (text generation), each at a fixed checkpoint revision
(MM~\mm{25.210}).

\begin{xltabular}{\linewidth}{@{}lXX@{}}
\caption{What the project supplies at each layer of the native path, and what is still Apple's or outside the project (MM~\mm{25.210.9}).}\label{tab:native-path-ownership}\\
\toprule
Layer & This project & Still Apple's, or outside \\
\midrule
\endfirsthead
\tablecontinued{3}
\toprule
Layer & This project & Still Apple's, or outside \\
\midrule
\endhead
model & graph import, checkpoint conversion, lowering, memory plan & the checkpoints and tokenizers; NumPy and PyTorch for host input and
validation \\*
GPU program & every instruction byte, from the compiler of \cref{ch:compiling} & encoding rules recovered with Apple's decoder \\*
execution & the host worker: uploads, resource and command state, Submit, wait,
checks & \code{IOGPU.framework}, IOKit, the AGX kernel driver, the firmware, the
hardware \\*
\bottomrule
\end{xltabular}

Both models pass through the same pipeline (MM~\mm{25.210.3}).

\begin{enumerate}
\def\labelenumi{\arabic{enumi}.}
\item \code{agxforge/g17/inferencegraph.py} reads the checkpoint's tensors into a typed graph and refuses holes,
  overlaps and mismatched shapes.
\item \code{tools/g17inferencelower.py} compiles each node with the compiler of \cref{ch:compiling}. MiniLM lowers to
  170 stages built from 16 distinct bodies, and Qwen to 748 prefill and 916 decode stages built from 28.
\item A memory planner assigns arenas and gives each slot 256-byte guards.
\item \code{tools/g17inferencebundle.py} places the schedule into a measured resource graph of 128 MiB for MiniLM or
  1 GiB for Qwen, and a resident worker runs it with at most three bindings and one Submit per stage.
\end{enumerate}

For each stage the runner places the stage's program in allocation 0 and writes the stage's launch words (the program
counter, the shader state, the grid, the threads per threadgroup and up to three buffer pointers). It then builds one
count-one Submit record with fresh callback blocks, calls IOGPU and waits up to five seconds for both callbacks.
\Cref{ch:submission-below-metal} documents these launch words and Submit records; the runtime refuses any value
outside the envelope measured there.

While the model runs, each stage reply carries bits recording that the stage completed, that the weights, guards,
code and uploaded source are unchanged, and that AGXMetal is still absent. The runner checks the guards of the live
regions after each stage, and any change in the GPU's recovery or event counters fails the run.

Correctness is checked offline on the captured outputs. Matrix operands travel as half with fp32 accumulation, and the
checks are against that domain rather than original-fp32 arithmetic. The bounds below were registered before
the first dispatch and never widened (MM~\mm{25.210.2}).

\begin{enumerate}
\def\labelenumi{\arabic{enumi}.}
\item The CPU recomputes every stage from the values its GPU predecessor produced, so a stage is checked against
  its actual input. The bound is \code{2e-5 * (1 + abs(reference))}, and packs, gathers,
  bias and residual adds must be bit-exact. All 680 MiniLM stages and 9,984 Qwen stage checks pass; drift across the
  model is outside what this level sees.
\item Application outputs are checked against the original checkpoint evaluated in FP64. MiniLM's final embedding
  must be within a maximum error of 0.005 and a cosine of at least 0.999; it measured 1.29e-4, with cosine at least
  0.9999997786. Against an FP64 model of the same half-transport policy, the bound \code{2e-4 * (1 + abs(reference))}
  also passes (maximum 8.21e-5). Each Qwen last-row logit vector, with the native tokens teacher-forced, must be within
  \code{0.05 + 0.003 * abs(reference)}. All 37 vectors fall inside, and all 37 agree on the top token.
\end{enumerate}

One auxiliary control fails. Against an FP64 model of the same half-projection policy, with the bound
\code{0.025 + 0.002 * abs(reference)}, Qwen's prefix logits reach a maximum error of 0.097 and 10,658 elements fall
outside. A replay of the failing prefix traces the excess to a half-rounding boundary in layer 0.

The native path is slower than the matched Metal blocks and MPS (MM~\mm{25.210.5}). The native MiniLM FFN block takes
2.427~ms against 1.797~ms for its matched Metal block, and the attention block 3.815 against 1.962~ms. In the paired
comparison with Transformers on MPS, native Qwen decode ran 3.26–3.33 tok/s against MPS's 22.6–49.1; the lowest MPS
figure is its first request, which includes first-use effects, and its warm requests ran 48–49. Later changes,
with no MPS run alongside, raised native decode to 5.15–5.32 tok/s in alternation with a build that kept the old guard
check, and to 4.62 tok/s in a separate replay. A full-model Metal control written in this repository runs the same 28
programs; its decode frame is 230~ms, of which 33~ms is reported GPU time.

The native path's per-token time has not been divided among its 916 Submits, guard scans and padded transfers.
\Cref{sec:completion-callbacks-their} gives the FFN block's interval from Submit to completed callback, and replacing
the host's scalar guard check with a vectorized one made Qwen decode 1.30–1.31× faster.

\section{Comparison of the two paths}\label{ch:comparing-two-paths}

The native path has not been shown to be faster than the path through Metal (\cref{tab:paths-compared}).

\begin{xltabular}{\linewidth}{@{}>{\hsize=0.55\hsize}XX>{\hsize=1.39\hsize}X@{}}
\caption{The path through Metal and the native path below it, side by side.}\label{tab:paths-compared}\\
\toprule
& through Metal & below Metal \\
\midrule
\endfirsthead
\tablecontinued{3}
\toprule
& through Metal & below Metal \\
\midrule
\endhead
models & InternLM2.5-1.8B, Qwen3-0.6B, Qwen3-8B & MiniLM-L6, Qwen2.5-0.5B \\*
pipeline configuration & Metal derives it from the image's metadata & the runner writes it, from measured classes \\*
speed & decode at 1.14–1.22× the rate of mlx-lm's \code{generate\_step} (MM~\mm{25.190}) & slower than matched Metal blocks and MPS; no advantage either way
against the full-model Metal control \\*
not established & why Apple's code for the same kernels is faster & a speedup from bypassing Metal (the two paths differ in submission, guards, padding and
transfers); behavior across OS updates (the IOGPU interfaces are private); reproduction outside this project; an
experiment that only this path can run (MM~\mm{25.210.9}) \\*
\bottomrule
\end{xltabular}
